\section{Lifetime checkpoints and reward selection}
\label{sec:method}

Choose $0<\rho<1$ before observing rewards. For arm $i$, set
\begin{equation}
 q_i(\rho)=\left\lceil\frac2{\rho^2}
                  \log\frac{4i(i+1)}\delta\right\rceil.
 \label{eq:checkpoint}
\end{equation}
Before this count is reached, its reporting interval is $[0,1]$. Afterwards,
the interval is permanently
\begin{equation}
 [\ell_i,u_i]=[\widehat\mu_{i,q_i}-\rho/2,
                 \widehat\mu_{i,q_i}+\rho/2]\cap[0,1].
 \label{eq:fixed_interval}
\end{equation}
The interval uses a fixed potential prefix, even though its completion time is
adaptive. Later observations continue to update the reward estimate. One can
also prescribe different precisions $\rho_i$ before observing the corresponding
arms; selecting a favorable checkpoint from their realized data is not covered.

CP-UCB-AA uses the following physical-clock rule.
\begin{enumerate}
\item Admit current arrivals and append their checkpoints to a queue in arrival
order. Initialize their counts to zero.
\item At an even round, if a checkpoint is incomplete, sample the first
incomplete arm in the queue. These are designated service pulls.
\item At every other round maximize, over all $i\in A_t$,
\begin{equation}
 J_i(t)=\widehat\mu_{i,N_i(t-1)}+
             \sqrt{\frac{3\log(2t)}{N_i(t-1)}}.
 \label{eq:ucb}
\end{equation}
Use $+\infty$ for a zero count and break ties in arrival order.
\item Update the chosen arm's lifetime count and reward sum. If the count first
equals $q_i$, construct \eqref{eq:fixed_interval}. Every observation contributes
to checkpoint completion, including observations from \eqref{eq:ucb}.
\item Set $\underline B_t=\max_{i\in A_t}\ell_i$ and
$\overline B_t=\max_{i\in A_t}u_i$, and report
\begin{equation}
 C_t=\{i:\ell_i\ge\overline B_t-\varepsilon\},\qquad
 D_t=\{i:u_i<\underline B_t-\varepsilon\},\qquad
 V_t=A_t\setminus(C_t\cup D_t).
 \label{eq:reports}
\end{equation}
\end{enumerate}
The intervals are nonempty on every sample path, so $C_t$ and $D_t$ cannot
overlap. All arrived arms remain eligible for \eqref{eq:ucb}, including those
in $D_t$. The method stores a count, sum, and reporting interval for every arm;
it does not claim to reduce total arm storage. A direct index or report update
costs $O(K_t)$ operations. The executable fixed-precision reference is supplied
with the numerical code.

\begin{proposition}[Simultaneous reporting and unconditional reward cost]
\label{prop:reporting}
On the full bounded model, CP-UCB-AA satisfies
$\Prob(E_{\rm lab})\ge1-\delta/2$. If $\mathcal I_T$ is its set of designated
service rounds, then at every deterministic $T$,
\begin{equation}
 \E R_T\le\min\left\{T,K_T+4\sqrt{3K_TT\log(2T)}+1+
                \E\sum_{t\in\mathcal I_T}\Delta_{I_t}(t)\right\}.
 \label{eq:reward_bound}
\end{equation}
Pathwise, $|\mathcal I_T|\le\sum_{i\le K_T}q_i$.
\end{proposition}

Hoeffding's inequality bounds the failure at checkpoint $i$ by
$\delta/\{2i(i+1)\}$. Summing over potential arm indices gives one coverage
event of probability at least $1-\delta/2$, on which every report is correct.
For \eqref{eq:reward_bound}, a separate, summable concentration argument at
physical time $t$ controls ordinary UCB pulls. Each arm and pre-pull count
is charged at most once; inserted service observations only skip counts.
This proof holds on all reporting paths and therefore incurs no omitted
$\delta T$ term. \suppproof

The construction retains the screening motivation of UCB-AA but changes its
algorithmic clock and its statistical task. It has no prescribed terminal round,
no delayed admission, and no restart that is completed artificially for analysis.
The guarantees above concern this construction, not the original round-based
elimination algorithm.
